\chapter{Fazit und Ausblick}
\label{chap:chapter-8}

Diese Projektarbeit hatte zum Ziel, eine interaktive grafische Benutzeroberfläche zu entwickeln, die die Bildklassifikation verständlicher macht. Durch das Tool \textit{Visual Explainable AI} sollte Studierenden der Einstieg in das Fachgebiet \textit{Explainable Artificial Intelligence} erleichtert werden.\\

Zu diesem Zweck wurden Anforderungen anhand User Stories und engen Gesprächen mit dem Hauptverantwortlichen erhoben und begründet analysiert. Daraus sind unter Anwendung der Moscow-Methode fünf Must-Have-, eine Should-Have-, eine Could-Have- und sechs Wont-Have-Anforderungen entstanden, die zur Erstellung des Softwarekonzepts verwendet wurden, das durch die statischen und dynamischen Modelle beschrieben wird.\\

So wurde das Tool \texttt{Visual Explainable AI} implementiert. \texttt{Visual Explainable AI} verwendet Technologien wie Python, FastAPI, PyTorch in Form einer Webapplikation zur Erzeugung von Vorhersagen (EfficientNetV2 und ResNet) und Erklärungen (Grad-CAM und LIME) für die Bildklassifikation.  \\ 

Die Webanwendung setzt die definierten Muss-Have-, Should-Have- und Could-Have-Anforderungen erfolgreich um: Sie verfügt über eine grafische Benutzeroberfläche, die es den Nutzern ermöglicht, Bilder hochzuladen, Vorhersagen zu erhalten und die Erklärungen für diese Vorhersagen zu visualisieren. Darüber hinaus ist das Tool wartbar und durch die MVP-Architektur leicht erweiterbar um weitere Neuronale Netze sowie Erklärungsverfahren, die sowohl textbasierte als auch grafische Erklärungen bieten.\\

Zur Evaluierung des Tools wurde eine Fallstudie durchgeführt, die den Einfluss unseres Tools auf Studierende und deren Verständnis der Grenzen von KI-Modellen, sowie die Nachvollziehbarkeit und den Einfluss von XAI Erklärungen untersucht. Die Ergebnisse der Fallstudie deuten darauf hin, dass die Mehrheit der Teilnehmer ein besseres Verständnis für die Bedeutung von XAI-Erklärungen durch die Anwendung des Tools entwickelt haben könnte. Aufgrund der geringen Teilnehmerzahl können jedoch keine endgültigen Schlussfolgerungen gezogen werden. \\

Alles in allem wurden die Ziele dieser Projektarbeit erreicht, da das Tool \texttt{Visual Explainable AI} erfolgreich implementiert. Die Ergebnisse der Evaluation der Auswirkungen des Tools auf das Verständnis von KI-Modellen durch Studierende sind vielversprechend, jedoch sind weitere Untersuchungen mit einer größeren Teilnehmerzahl sind erforderlich, um diese Ergebnisse zu bestätigen.\\